\documentclass[10pt,a4paper]{amsart}
\usepackage[margin=19mm]{geometry}
\usepackage{amsmath,amssymb,mathtools}
\usepackage[scr=rsfs,cal=euler]{mathalfa}
\usepackage{xfrac}
\usepackage[unicode,hidelinks,bookmarks=false]{hyperref}
\newcommand{\ie}{i.e.\ }
\newcommand{\cf}{cf.\ }
\newcommand{\resp}{respectively\ }
\newcommand{\Subsection}{\subsection}
\providecommand{\nobreakdash}{\nobreak\hbox{-}}
\setlength{\emergencystretch}{2em}
\multlinegap0pt
\usepackage{amssymb}
\usepackage{mathtools}
\usepackage[shortlabels]{enumitem}
\usepackage{xcolor}
\usepackage{booktabs}
\usepackage{multirow}
\usepackage{caption}
\newtheorem{prop}{Proposition}
\newcommand{\RR}{{\mathbb{R}}}
\newcommand{\dom}{\mathop{\rm dom}}
\title{{\Large Nesterov's Accelerated Gradient for Unbounded Convex Functions\\Finds the Minimum-Norm Point in the Dual Space}}
\author{Keiya Sakabe}
\date{Source: arXiv:2602.08618v1 (CC BY 4.0)}
\begin{document}
\maketitle

\noindent\emph{Source and licence.} {\Large Nesterov's Accelerated Gradient for Unbounded Convex Functions\\Finds the Minimum-Norm Point in the Dual Space}. Version arXiv:2602.08618v1 (CC BY 4.0), distributed by arXiv under the Creative Commons Attribution 4.0 International licence, http://creativecommons.org/licenses/by/4.0/, as recorded in the arXiv licence metadata for this exact version and shown on the abstract page https://arxiv.org/abs/2602.08618v1. Full licence notice: \texttt{licenses/arxiv-2602.08618-CC-BY-4.0.txt}.\par\smallskip

\noindent\textit{Editorial scope.} Counted unit: the article's prop prop:boundeddual (source0.tex lines 622--629). The article's complete original proof of it is delivered with it (source0.tex lines 630--638, one \texttt{proof} environment of 9 lines). No mathematics is written, repaired or re-derived: the statement and the proof are the article's own lines, in the article's own notation, and no retained line is substituted or re-flowed.  No further source-local context is needed: every cross-reference of every retained line resolves inside the two delivered spans.  Nothing was omitted from any retained span, so the package carries no omission record.  No retained line contains a citation, so the wrapper carries no bibliography and no bibliography entry.  No retained line contains a figure environment, an image-inclusion command or drawing code, so no image is delivered and the package declares no asset dependency.  Editorial formatting only, in addition: an AMS \texttt{amsart} preamble that restates the article's own declarations and macros used by the retained lines, namely the environments \texttt{prop} on separate counters, together with the standard packages those lines need.  The retained environments print in this document as Proposition 1 in order of appearance, whereas the article prints the same items with the numbering of its own full document.  The complete native source of the article as distributed in the arXiv e-print for this exact version is bundled unchanged as \texttt{sources/194259/source0.tex}.
\begin{prop}\label{prop:boundeddual}
    Suppose $f^\ast \le M$ for some $M \in \RR$.
    Then, $\dom f^\ast$ has a minimum-norm point $p^\star$.
    Moreover, for any $x_0 \in \RR^n$, it holds that
    \[
        D_f(x_0, p^\star) \le M + f(x_0) + \|x_0\|\|\nabla f(x_0)\|.
    \]
\end{prop}

\begin{proof}
    Since $f^\ast$ is lower semicontinuous, for any $p \in \overline{\dom f^\ast}$, it holds that $f^\ast(p) \le M < \infty$.
    Therefore, $\dom f^\ast$ is closed; in particular, $\dom f^\ast$ has a minimum-norm point $p^\star$.
    Then, for any $x_0 \in \RR^n$,
    \[
        D_f(x_0, p^\star) = f(x_0) + f^\ast(p^\star) - \langle p^\star, x_0\rangle \le M + f(x_0) + \|x_0\|\|\nabla f(x_0)\|,
    \]
    where we used $\|p^\star\| \le \|\nabla f(x_0)\|$ for the inequality.
\end{proof}

\end{document}
